%!TEX program = lualatex
\documentclass[10pt, openany]{book}

\usepackage{comment}
\usepackage{silence}
\WarningFilter{latex}{Command \underbar has changed}
\WarningFilter{latex}{Command \underline has changed}

% ============== 考虑出血的版式 ===================================
\usepackage[
    paperwidth=6in,
    paperheight=9in,
    top=0.8in,
    bottom=0.8in,
    inner=0.75in,
    outer=0.6in,
    bindingoffset=0.25in,
    footskip=0.4in
]{geometry}

% ============== 语言和字体 ===================================
\usepackage[x-4]{pdfx} 
\usepackage{fontspec}
\usepackage[provide=*, chinese]{babel}
\usepackage{sectsty}

\babelfont{rm}[
    Path = ./font/,
    Renderer = HarfBuzz,
    BoldFont = {GFSDidot-Regular.ttf},
    BoldFeatures = {RawFeature={embolden=3.0}},
    ItalicFont = {timesi.ttf}
]{GFSDidot-Regular.ttf}

\babelfont{sf}[Path = ./font/, Renderer = HarfBuzz]{GFSDidot-Regular.ttf}
\babelfont{tt}[Path = ./font/, Renderer = HarfBuzz]{NotoSerifSC-Regular.ttf}

\babelfont[chinese]{rm}[
    Path = ./font/,
    Renderer = HarfBuzz,
    BoldFont = {NotoSerifSC-ExtraBold.ttf},
    ItalicFont = {StKaiTi.ttf}
]{NotoSerifSC-Regular.ttf}

\allsectionsfont{\rmfamily\bfseries}

\usepackage{amssymb}
\usepackage{enumitem}
\usepackage{titlesec}
\usepackage{quoting}
\usepackage{setspace}
\onehalfspacing
\usepackage{tikz}
\usepackage{fancyhdr}
\usepackage{amsmath, amsthm}
\usepackage{bm} 
\usepackage{libertine} 
\usepackage[libertine]{newtxmath} 
\usepackage{tikz-cd} 
\usepackage{nicematrix} 
\usepackage{adjustbox} 
\usepackage{framed} 
\usepackage{scrextend}
\changefontsizes{9.5pt} 


\NiceMatrixOptions{
  custom-line = {
    letter = : ,
    command = hdashline ,
    tikz = { dashed }
  }
}

\setlength{\parskip}{1em} 

\pagestyle{fancy}
\fancyhf{}
\fancyhead[LE,RO]{\thepage}
\fancyhead[RE]{\textbf{西 \quad 窗 \quad 烛}}
\fancyhead[LO]{\textit{The Candle by the West Window}}
\renewcommand{\headrulewidth}{0.4pt}

\begin{document}

\frontmatter
\title{\textbf{量子力学讲义}\\ \large 概率的复算子化}
\author{Lao Chen}
\date{\today}
\maketitle
\tableofcontents

\mainmatter

\chapter{Hamilton力学}

简单起见讨论一维空间的力学，
能量$e$的变化即为功$\text d e$，
力$F$按照位移$\text d s$的做功：

$$
\text d e = F \text d s
$$

牛顿定律，
把力$F$视为动量$p=mv$随时间的变化率，
动量视为力随时间的积分：

$
F = ma
= m\frac{\text dv}{\text dt}
= \frac{\text dp}{\text dt}
$

动量视为动能$T$随时间的变化率，
动能视为动量随时间的积分：

$$
\frac{\text dT}{\text dt}
=\frac{m\text d v^2}{2 \text dt}
=mv = p
$$

用Legendre变换，
把动能写为动量的函数：

$$
T(p) = \frac{p^2}{2m}
$$

令$F(x)$是一种由位置$x$决定的力，
沿着位置积分得到：

$$
V = -\int F(x) \text d x
$$

它反映了克服力的方向做功所累积的能量，
即势能，
反过来：

$$
F(x) = -\frac{\text d V}{\text d x}
$$

总能量构成Hamilton量：

$$
H(x,p) = \frac{p^2}{2m} + V(x)
$$

有

$$
\frac{\partial H}{\partial x}
=\frac{\partial V}{\partial x}
= -F, \quad 
\frac{\partial H}{\partial p}
=\frac{p}{m}
= v
$$

\chapter{量子力学公理}

\section{公理1：量子Hilbert空间}

完备的内积空间称为Hilbert空间。
系统的状态由Hilbert空间$H$中的单位向量$\psi$表示。
若$\exists c\in\mathbb C$使得$\psi_1 = c\psi_2$，
则构成等价关系$\psi_1 \sim \psi_2$，
二者代表同样的状态。
这样的$H$称为量子Hilbert空间。

\section{公理2：自伴算子}

经典相空间的函数$f$，
关联到量子Hilbert空间$H$上的自伴算子$\hat f$。

\section{公理3：测量}

如果系统处于状态$\psi \in H$，
则可观察量$f$的测量概率分布满足：

$$
E(f^m) = \langle \psi,(\hat f)^m\psi \rangle
$$


\chapter{有限概率分布}

有限集合$X\in\textbf{Set}$
对应一个$|X|$-维线性空间$V\in\textbf{FdVct}$，
令有正交归一基$E = \{e_i \in V\}$，
其中每个基向量$e_i$张成一个$1$-维子空间$E_i$，
有：

$$
V = \bigoplus_i E_i
$$

概率的区间为闭区间$I = [0,1]$，
在$X$上定义有限概率分布
$p: X \to I :: x_i \mapsto p_i = p(x_i)$，
满足$\sum_i p_i = 1$。

记$V$上自伴算子的全体为$\text{Herm}(V)$。
任何自伴算子$\forall A\in\text{Herm}(V)$都有实特征值，
即$\sigma(A) = \{ \lambda_i \} \subset \mathbb R$，
且适当选择基可以使
$A = \text{diag}(\lambda_i)$，
并使其对应酉算子
$U = \text e^{\text i A} = \text{diag}(\text e^{\text i \lambda_i})$。

以集合的第$i$个元素对应的概率作为特征值$\lambda_i = p_i = p(x_i)$，
得到实对称矩阵
$A = \text{diag}(p_i)$。
概率的性质使得这个矩阵具有：

\begin{itemize}
  \item 半正定
  \item 归一化：$\text{tr} \ A = 1$
\end{itemize}

\chapter{实数化的概率}

集合$X$上有幂集$PX$，
幂集的子集为$X$的子集族，
若子集族$\Omega$构成$\sigma$-代数，
且在$\Omega$上定义了非负的测度$\mu$，
则$(X,\Omega,\mu)$构成测度空间。

如果$X$可以表示为可数个具有有限测度的集合的并，
称为$\sigma$-有限的。

\textbf{Radon-Nikodym定理}：
若$(X,\Omega)$上有两个$\sigma$-有限测度$\mu$和$\nu$，
$\mu$对于$\nu$绝对连续，
则存在$X$上的非负可测函数$\rho$使得：

$$
\mu(E) = \int_E \rho \text d \nu
$$

对所有$E\in\Omega$成立，
$\rho$称为$\mu$对于$\nu$的密度。

\chapter{复数化的概率}

\section{作为实线性算子的复数}

实数的乘法$rv$，
视左边的实数$r$为线性算子，
视右边的实数为$1$维护向量输入，
用矩阵对列向量作用的方式写为：

$$
\begin{aligned}
[r]: \mathbb R_1 &\to \mathbb R_1 \\
[v] &\mapsto [r] [v] = [rv] \\  
\end{aligned}
$$

从$1$维实线性空间$\mathbb R_1$到其自身的实线性算子全体，
构成实线性空间：

$$
\textbf{End}(\mathbb R_1) = \textbf{Vct}(\mathbb R_1,\mathbb R_1)
$$

作为矩阵的$[r]\in\textbf{End}(\mathbb R_1)$，
和作为实数的$r\in\mathbb R$是一一对应的。
若有两个算子$p,q\in \textbf{End}(\mathbb R_1)$，
其复合作用是矩阵乘法，
相当于实数乘法：

$$
\begin{aligned}
[p] \circ [q]: \mathbb R_1 &\to \mathbb R_1 \\
[v] &\mapsto ([p] \circ [q]) [v] = [pqv] \\  
\end{aligned}
$$

有了乘法就可以考虑幂和开根号。
在线性映射的上下文中，
实数$1$相当于恒等线性映射$[1]$，
而实数$-1$相当于镜像映射$[-1]$，
于是我们需要考虑是否存在线性算子$[i]\in\textbf{End}(\mathbb R_1)$，
使得$[i] \circ [i] = [-1]$呢？
这就是$i = \sqrt{-1}$这个开根号问题在实线性算子角度的描述，
答案是否定的。
进而放在二维考虑类似问题，
有$2\times 2$维的矩阵$T\in\textbf{End}(\mathbb R_2)$作为线性映射：

$$
\begin{aligned}
A: \mathbb R_2 &\to \mathbb R_2 \\
\begin{bmatrix}v_1\\v_2\end{bmatrix} &\mapsto 
\begin{bmatrix}a & b\\c & d\end{bmatrix} \begin{bmatrix}v_1\\v_2\end{bmatrix} \\  
\end{aligned}
$$

进而我们继续追问，
在二维情况下是否可以把镜像映射用开平方根的问题，即：

$$
\begin{aligned}
\begin{bmatrix}a & b\\c & d\end{bmatrix}
\begin{bmatrix}a & b\\c & d\end{bmatrix}
= \begin{bmatrix}-1 & 0\\0 & -1\end{bmatrix}
\end{aligned}
$$

在二维的情况，
镜像相当于角度为$\pi$的旋转，
而旋转是可以矩阵表示的线性算子：

$$
R_\theta = \begin{bmatrix}
\cos \theta & \sin \theta \\
-\sin \theta & \cos \theta \\
\end{bmatrix}
$$

显然对镜像的根号即为$\theta = \pi/2$的旋转矩阵，
它就是虚数单位：

$$
\text i
=\begin{bmatrix}
\cos \theta & \sin \theta \\
-\sin \theta & \cos \theta \\
\end{bmatrix}
=\begin{bmatrix}
0 & 1 \\
-1 & 0 \\
\end{bmatrix}
$$

对于Euler公式有了二维实矩阵表示：

$$
e^{\text i\theta} = \cos \theta + \text i\sin \theta 
=\cos \theta\begin{bmatrix}1&0\\0&1\end{bmatrix}
+\sin \theta\begin{bmatrix}0&1\\-1&0\end{bmatrix}
=\begin{bmatrix}
\cos \theta & \sin \theta \\
-\sin \theta & \cos \theta \\
\end{bmatrix}
$$

这是单位圆上的复数，
进而推广到任意复数有：

$$
z = x + \text iy =
\begin{bmatrix}
x & y \\ -y & x
\end{bmatrix}
$$

\section{共轭}

再看一个经典公式$\cos^2\theta + \sin^2\theta = 1$的矩阵表示：

$$
\begin{aligned}
1 &= e^0 = e^{\text i\theta}e^{-\text i\theta} 
= (\cos \theta + \text i\sin \theta)(\cos \theta - \text i\sin \theta) \\
&=\begin{bmatrix}
\cos \theta & \sin \theta \\
-\sin \theta & \cos \theta \\
\end{bmatrix}
\begin{bmatrix}
\cos \theta & -\sin \theta \\
\sin \theta & \cos \theta \\
\end{bmatrix}
=\begin{bmatrix}
1 & 0 \\
0 & 1 \\
\end{bmatrix}
\end{aligned}
$$

它反映了单位圆上的复数，
其共轭相当于反向旋转。
进而考虑任意复数的共轭：

$$
z = x + \text iy =
\begin{bmatrix}
x & y \\ -y & x
\end{bmatrix}, \quad 
\overline z = x - \text iy =
\begin{bmatrix}
x & -y \\ y & x
\end{bmatrix}, \quad 
$$

其乘积为：

$$
\begin{aligned}
z \overline z
&= (x + \text iy)(x - \text iy)
= x^2 + y^2 = |z|^2 \\
&=\begin{bmatrix}
x & y \\ -y & x
\end{bmatrix}
\begin{bmatrix}
x & -y \\ y & x
\end{bmatrix}
=\begin{bmatrix}
x^2 + y^2 & 0 \\ 0 & x^2 + y^2
\end{bmatrix}
\end{aligned}
$$

\section{实数概率}

在单位圆盘$D = \{z \in \mathbb C \ | \ |z| \leq 1\}$内部的任意复数$z$，
其模平方$|z|^2 = z \overline z$的取值范围是$I = [0,1]$。
这个实数中的闭区间，
正是概率的取值范围。
这样，任何概率$p$都存在$p = |z|^2$的表示。
并且，它是旋转不变的。
换句话说，
一个实数概率$p\in I$，
和复平面单位圆盘$D$中的一个半径为$\sqrt{p}$的圆一一对应。
若两个复数相差一个旋转，
则$|z_1| = |z_2|$，
构成等价关系$z_1 \sim z_2$，
于是概率的单位闭区间可以描述为商集：

$$
I = [0,1] = D/\sim
$$

在满足某些概率条件时，
其上有可加性：

\begin{equation}
p_1 + p_2 \in I = [0,1]
\end{equation}
\label{eq:real-prob-sum}

直接在实数上相加是概率论讨论的问题。
而放到复平面上，
对应了半径为$\sqrt{p_1}$和$\sqrt{p_2}$的两个圆的某种加法，
这似乎比实数概率的直接相加还缺少直观意义。
如果考虑两个圆上的点的复数加法则有趣的多。

\section{复数化的概率}

假设$p_1 = |z_1|^2$和$p_2 = |z_2|^2$，
且$|z_1| \geq |z_2|$，
复数加法相当于线性算子的求和：

$$
z_1 + z_2
=\begin{bmatrix}x_1 & y_1 \\ -y_1 & x_1\end{bmatrix}
+ \begin{bmatrix}x_2 & y_2 \\ -y_2 & x_2\end{bmatrix}
= \begin{bmatrix}x_1 + x_2 & y_1 + y_2 \\ -y_1 - y_2 & x_1 + x_2 \end{bmatrix}
$$

其模平方为：

\begin{equation}
\begin{aligned}
|z_1 + z_2|^2
&= (x_1 + x_2)^2 + (y_1 + y_2)^2 \\
&= x_1^2 + 2x_1x_2 + x_2^2 + y_1^2 + 2y_1y_2 + y_2^2 \\
&= x_1^2 + y_1^2 + x_2^2 + y_2^2 + 2x_1x_2 + 2y_1y_2 \\
&= |z_1|^2 + |z_2|^2 + 2(x_1x_2 + y_1y_2) \\
&= p_1 + p_2 + 2(x_1x_2 + y_1y_2) \\
\end{aligned}
\end{equation}
\label{eq:complex-prob-sum}

对比在商集$I = D/\sim$中的实数加法(\ref{eq:real-prob-sum})，
这里在复平面上的运算要灵活的多。
在上式里我们看到$p_1 + p_2$出现了，
但最终的概率受到$2(x_1x_2 + y_1y_2)$项的修正，
它反映了两个复数之间的相位差夹角，
使得最终的概率取值范围为：

$$
|z_1|^2 - |z_1|^2 \leq |z_1 + z_2|^2 \leq |z_1|^2 + |z_1|^2
$$

这反映出，
尽管两个复数分别通过模平方携带了各自的概率信息，
但两者做复数加法时，
可能会互相抵消一部分，
使得求和后对应的概率介于以上范围。
于是我们把实数概率求和发展为某种可以互相抵消一部分概率的求和。

\chapter{泛函分析}


\section{交换代数}

集合范畴$\textbf{Set}$中，
用反变Hom函子$h^Z$作用于集合$X$得到函数态射集
$Z^X = h^Z X = \textbf{Set}(X,Z)$。

若加强为Abel群$A\in\textbf{Ab}$，
则$A^X = \textbf{Set}(X,A)\in\textbf{Ab}$按照$R$的逐点加法运算构成Abel群或者$\mathbb Z$-模。

若加强为交换环$R\in\textbf{CRing}$，
则$R^X = \textbf{Set}(X,R)\in\textbf{CRing}$按照$R$的逐点乘法运算构成交换环。
这样可以讨论其中的理想。

继续加强为域$k$，
则$k^X = \textbf{Set}(X,k)\in\textbf{Vct}_k$按照$k$的逐点数乘运算构成$k$-线性空间或者$k$-模。

再加上交换的乘法，
使其构成含幺交换结合代数。

\section{算子代数}

$k$-线性空间$V\in\textbf{Vct}_k$上考虑线性算子的全体
$\textbf{End} \ V = \textbf{Vct}_k(V,V)$。
为算子$A\in\textbf{End} \ V$定义算子范数：

$$
\| A \| = \sup \frac{\| A\psi \|}{\| \psi \|},\quad 
\psi \in V
$$

若算子范数$\| A \|$有限称$A$为有界算子。
$\mathcal B(V) \subseteq \textbf{End} \ V$是有界算子的集合。
对于$1$维以上的空间，
它构成非交换代数。
对于可分复Hilbert空间$H$，
$\mathcal B(H)$满足：

$$
\| A B \| \leq \| A \| \| A \|
$$

这使得$\mathcal B(H)$构成Banach代数。
复共轭使其构成C*-代数。

\section{连续函数代数}

令$X\in\textbf{LCT}$为局部紧拓扑空间。
连续函数的全体$\mathbb C^X = \textbf{Top}(X,\mathbb C)$构成交换C*-代数。
其幺元$\textbf 1$就是在$X$上取值恒为$1$的唯一的函数。

令$C_\infty(X) \subseteq \mathbb C^X = \textbf{Top}(X,\mathbb C)$为在无穷远处消失的连续函数全体。
$C_\infty(X)$是交换C*-代数。
只有当$X\in\textbf{CTop}$为紧拓扑空间时，
$\textbf 1 \in C_\infty(X)$，
于是代数的含幺性和拓扑的紧性有了对应。
含幺时将这个交换函数代数记为$C(X)$。

\section{Gelfand拓扑}

Gelfand-Mazur定理告诉我们，
含幺交换Banach代数$A \in \textbf{UCB}$对其极大理想$\mathfrak m\in\text{MaxSpec} \ A$做商代数有：

$$
A/\mathfrak m \simeq \mathbb C
$$

$a \in A$作为代表元，
代表了商代数中的元，
根据上式它可以直接视为复数
$a + \mathfrak m \in \mathbb C$。
换句话说我们得到一个二元复值函数：

$$
\begin{aligned}
A \times \text{MaxSpec} \ A &\to A/\mathfrak m \simeq \mathbb C \\
(a,\mathfrak m) &\mapsto a + \mathfrak m 
\end{aligned}
$$

固定一个变元$a \in A$得到极大谱上的复值函数：

$$
\begin{aligned}
\varphi_a: \text{MaxSpec} \ A &\to \mathbb C \\
\mathfrak m &\mapsto \varphi_a(\mathfrak m) = a + \mathfrak m 
\end{aligned}
$$

构造这个函数的过程是代数同态，
称为$A$的Gelfand表示：

$$
\begin{aligned}
A &\to \mathbb C^{\text{MaxSpec}} \\
a &\mapsto \varphi_a 
\end{aligned}
$$



\chapter{复线性算子}

\section{概率的保持}

在$n$元有限集上有概率分布$\{p_i\}_{1 \leq i \leq n}$，
满足$\sum p_i = 1$。
它可以用线性算子描述为：

$$
A = \text{diag}(p_1,\cdots,p_n) 
= \begin{bmatrix}
p_1 & & \\ & \ddots & \\ & & p_n
\end{bmatrix}
$$

对角矩阵$A$满足$\text{trace} \ A = 1$，
它对应了概率归一化的要求。

\section{自伴算子}

\section{酉算子}

\chapter{算子代数}

复可分Hilbert空间$H$上的线性算子$A$，
其算子范数：

$$
\|A\| = \text{sup}_{\psi \in H - \{0\}} \frac{\|A \psi\|}{\|\psi\|}
$$

都是有限的，
则称其为有界算子。
所有这些$H$上的有界算子，
按上述算子范数构成Banach空间，
记为$\mathcal B(H)$。
对于$\forall A,B \in \mathcal B(H)$有：

$$
\|AB\| \leq \|A\|\cdot\|B\|
$$

使其构成Banach代数。

\textbf{Banach代数}：
复Bananch空间$A$是一个复结合代数，
$A$称为Banach代数，
若其代数乘法满足

$$
\|xy\| \leq \|x\|\cdot\|y\|
$$

Banach代数$\mathcal B(H)$由$H$上的有界线性算子构成，
算子$A$的伴随算子记为$A^*$。
以伴随运算定义共轭运算，
使得$\mathcal B(H)$构成C*-代数。
当$\dim H > 1$时，
它是非交换代数。

另一个C*-代数的例子：
令$\Omega$为局部紧空间，
$\Omega$上复值连续无穷远消失的函数全体，
记为$C_\infty(\Omega)$，
构成C*-代数。
它是交换代数。
$C_\infty(\Omega)$是含幺代数的充分必要条件是$\Omega$是紧的，
此时记为$C(\Omega)$。

\chapter{谱理论}

$\lambda x$视为用实数$\lambda$作用于输入$x$，
这是一维的情况。
从二维开始需要用矩阵$A$作用于输入向量$x$得到$Ax$。
如果希望矩阵的作用方式$Ax$相当于标量的缩放$\lambda x$：

$$
(A - \lambda I)x = 0
$$

相当于向量$x$落在矩阵$A - \lambda I$作为线性映射的核：

$$
x\in\text{Ker}(A - \lambda I) 
= \{x \ | \ (A - \lambda I)x = 0\}
$$

如果矩阵$A - \lambda I$是可逆的，
那么核退化为原点，
只存在唯一的平凡解$x=0$。
对于非退化的情况，
若存在$x$和$\lambda$令上式成立，
则称为特征向量、特征值。

看平面旋转矩阵的谱：

$$
R_\theta = \begin{bmatrix}
  \cos\theta & -\sin\theta \\
  \sin\theta & \cos\theta \\
\end{bmatrix}
$$

$$
\begin{aligned}
\det (R_\theta - \lambda I)
&=\det \begin{bmatrix}
  \cos\theta - \lambda & -\sin\theta \\
  \sin\theta & \cos\theta -\lambda \\
\end{bmatrix} \\
&=P(\lambda)
=\lambda^2 - 2\cos\theta \lambda + 1 \\
&=0
\end{aligned}
$$

得到解集也就是谱
$\sigma(R_\theta) = \{ e^{\text i\theta},e^{-\text i\theta} \}$。
旋转不改变向量长度。
$|R_\theta x| = |x|$
对所有向量$x$成立，
体现在
$|\lambda| = |e^{\pm\text i\theta}| = 1$。
保证了算子是\textbf{酉(unitary)}的。

$$
R_\theta x = \begin{bmatrix}
  \cos\theta & -\sin\theta \\
  \sin\theta & \cos\theta \\
\end{bmatrix}x
= \lambda x
= \begin{bmatrix}
  e^{\pm\text i\theta} & 0 \\
  0 & e^{\pm\text i\theta} \\
\end{bmatrix}x
$$

再看一个普通矩阵的谱：

$$
A = \begin{bmatrix}
  2 & 1 & 0 \\
  0 & 2 & 1 \\
  -8 & 0 & 2
\end{bmatrix}
$$

$$
\begin{aligned}
\det (A - \lambda I)
&=\det \begin{bmatrix}
  2 - \lambda & 1 & 0 \\
  0 & 2 - \lambda & 1 \\
  -8 & 0 & 2 - \lambda
\end{bmatrix} \\
&=P(\lambda)
=-\lambda^3 + 6\lambda^2 - 12\lambda
=-\lambda(\lambda^2 - 6\lambda + 12) \\
&=0
\end{aligned}
$$

得到解集也就是谱
$\sigma(A) = \{ 0, 3 + \text i\sqrt{3}, 3 - \text i\sqrt{3}\}$。

\chapter{Gelfand理论}

\section{交换代数}

令$A\in\textbf{CRing}$为交换环，
子集$\mathfrak a \subseteq A$称为理想，
若$A \mathfrak a \subseteq \mathfrak a$。
不被任何非$(1)$包含的理想$\mathfrak m$称为极大理想。

$\mathfrak p$是素理想 iff. $A/\mathfrak p$是整环。
$\mathfrak m$是极大理想 iff. $A/\mathfrak m$是域。

矩阵的特征值就是特征多项式环上的素谱。

整数环$\mathbb Z$的理想是任何整数$n$构成的$(n)$。
素数$p$构成素理想$\mathfrak p = (p)$。
素理想的全体为素谱$X = \text{Spec} \ \mathbb Z$，
将其视为空间来讨论，
由于$\mathbb Z$是整环，
除素数构成的素理想作为闭点外，
$(0)\in X$是泛点，
它的闭包是整个$\text{Spec} \ \mathbb Z$。
讨论作为空间的素谱上的$\mathbb Z$-值函数：

$$
\begin{aligned}
f: X &\to \mathbb Z \\
\mathfrak p = (p) &\mapsto f(\mathfrak p) = f \mod p \\
\end{aligned}
$$

函数$f$在点$\mathfrak p$处的取值，
为$f$在剩余域$\mathbb Z/p\mathbb Z$中的投影。

紧空间$X$上的复连续无穷远消失函数所构成的交换代数$C(X)$中，
给定$E \subseteq X$，
当$x\in E$时$f(x) = 0$的函数构成理想$\mathfrak a_E$。


\chapter{Topos理论}



\end{document}